% ============================================================
%  Modul 8 — Implementasi Multi-threading
% ============================================================

\chapter{Implementasi Multi-threading}
\label{chap:modul-08}

% --- 1. Tujuan dan Output Praktikum ---
\section{Tujuan dan Output Praktikum}

\subsection{Tujuan Praktikum}
Tujuan yang diharapkan dapat dicapai oleh mahasiswa adalah:
\begin{enumerate}
    \item Memahami arsitektur perangkat lunak dalam manajemen \textit{thread} (termasuk abstraksi \textit{Thread Control Block}) dan algoritma penjadwalan (\textit{scheduling}) pada tingkat inti sistem operasi (\textit{kernel-level}).
    \item Mempraktikkan implementasi logika \textit{Context Switching} untuk memfasilitasi eksekusi \textit{multi-threading} melalui penyimpanan dan pemuatan status register CPU (memanfaatkan manipulasi bahasa tingkat rendah/Assembly).
\end{enumerate}

\subsection{Output Praktikum}
Hasil luaran yang diekspektasikan adalah kelengkapan instrumen berikut:
\begin{itemize}
    \item Implementasi fungsional dari modifikasi berkas sumber (khususnya penyesuaian pada logika \texttt{thread.c}, \texttt{switch.asm}, dan \texttt{scheduler.h}) yang terbukti mampu mengompilasi rekayasa antrean \textit{thread}.
    \item Tangkapan layar (\textit{screenshot}) keluaran terminal dari emulator QEMU yang memvalidasi bahwa beberapa \textit{thread} (contoh: Thread A, B, dan C) sukses berjalan saling bergantian secara sinkron mengambil alih siklus CPU.
\end{itemize}

% --- 2. Dasar Teori ---
\newpage
\section{Dasar Teori}

\subsection{Pengenalan Thread dan Multi-threading}
Dalam sistem komputasi, \textit{thread} adalah unit instruksi terkecil yang dapat dijadwalkan oleh sistem operasi \cite{Silberschatz2018}. Jika sebuah proses diibaratkan sebagai program yang sedang berjalan dan memiliki memori sendiri, maka \textit{thread} adalah entitas eksekusi operasional di dalam proses tersebut. Lingkungan OS \textit{bare-metal} merintis kerangka \textit{multi-threading} agar dapat mensimulasikan penanganan banyak fungsi bahasa C secara bersamaan (konkurensi) \cite{pintos_stanford}. Secara konseptual, rutinitas pembuatan dan pengelolaan \textit{thread} dikonstruksi guna mendistribusikan beban kerja secara seimbang sehingga mendongkrak efisiensi utilitas prosesor.

\subsection{Konteks Eksekusi dan Context Switching}
Konteks eksekusi merujuk pada seluruh status kumpulan register dari prosesor (seperti EAX, EBX, EBP, dan ESP pada perangkat arsitektur x86) yang dikenakan oleh sebuah \textit{thread} pada momentum waktu tertentu \cite{bryant2015computer}. Ketika sistem operasi hendak memberhentikan satu \textit{thread} lalu beralih ke giliran antrean eksekusi milik \textit{thread} yang lain, OS wajib menukarkan keadaannya lewat taktik perpindahan \textit{Context Switching}.

Karena kerumitan modifikasi letak tumpukan memori (\textit{stack}) amat bergantung pada wujud arsitektur silikon sistem biner, prosedur krusial ini umumnya direpresentasikan murni melalui arahan instruksi pada level \textit{Assembly}. Operasi tersebut mutlak ditugaskan untuk mengekstrak simpanan (\textit{Push}) dari *thread* lama seraya mengangkat muatan register (\textit{Pop}) profil milik antrean *thread* pengganti.

\subsection{Thread Control Block / TCB}
Agar sistem pengontrol sanggup menjaga identitas serta alokasi titik perhentian terakhir masing-masing \textit{thread}, direkayasa sebuah skema pelacak terstruktur dinamakan \textit{Thread Control Block} (TCB) \cite{Silberschatz2018}. Menggunakan struktur data mandiri (seperti halnya tipe rekaman mendasar pada bahasa C), blok TCB secara dedikatif membungkus parameter-parameter memori berikut:
\begin{itemize}
    \item \textbf{Thread ID}: Angka unit unik pengenal subyek yang tak boleh identik.
    \item \textbf{Status Thread}: Peta kondisi berjalannya instruksi pada suatu \textit{thread} (contoh: status sedang fasif siap diterjunkan/\textit{Ready}, status mandek/\textit{Blocked}, atau sedang berjalan/\textit{Running}).
    \item \textbf{Stack Pointer (ESP)}: Penunjuk referensi baris alamat terakhir yang terekam pada tumpukan (\textit{stack}) ketika interupsi perpindahan prosesnya berlangsung.
\end{itemize}

\subsection{Algoritma Penjadwalan (Scheduling)}
Siklus perancangan eksekusi CPU dikoordinasi secara tersentralisasi oleh modul OS benama \textit{Scheduler} \cite{Silberschatz2018}. Esensi utama dari algoritma komponen ini bertugas menyeleksi manakah prioritas \textit{thread} yang menduduki gerbang \textit{Ready} pada antrean sistem, agar mendapat wewenangnya dialihkan ke tingkatan fase \textit{Running}. Konsep paradigma fundamental penyeleksian dijabarkan melalui dua cara:
\begin{enumerate}
    \item \textbf{Cooperative Scheduling}: Sebuah \textit{thread} secara sadar membiarkan algoritma lain ganti mengeksekusi CPU dengan inisiatifnya menyerah (seperti metode inkuiri \texttt{yield()} atau tertahan menunggu sumber daya antarmuka selesai).
    \item \textbf{Preemptive Scheduling}: Sistem operasi menyetel batasan pencacah peranti keras layaknya \textit{Timer Interrupt}, berfungsi untuk merampas status dominasi algoritma tanpa ampun dari instruksi \textit{thread} yang sedang menduduki memori apakala batasan waktunya telah dirasa mencukupi.
\end{enumerate}

Tabel \ref{tab:m08-dasar-teori} memberikan intisari terminologi vital terkait landasan teori \textit{multithreading}.

\begin{table}[!ht]
\caption{Tabel Ringkasan Dasar Teori Multi-threading}
\label{tab:m08-dasar-teori}
\centering
\begin{tabulary}{\textwidth}{|L|L|}
\hline
\textbf{Konsep} & \textbf{Penjelasan Singkat} \\ \hline
Thread & Unit eksekusi fungsional terkecil di dalam OS. \\ \hline
Context Switch & Tindakan rekam-simpan silang arsitektur register satu \textit{thread} dengan \textit{thread} cadangan lainnya. \\ \hline
TCB & Struktur blok tipe data spesifik pengingat status profil \textit{thread}. \\ \hline
Scheduler & Organ algoritma OS yang mendikte hak pembagian waktu kendali CPU atas \textit{thread}. \\ \hline
\end{tabulary}
\end{table}

% --- 3. Alat dan Bahan ---
\newpage
\section{Alat dan Bahan}
\label{sec:m08-alat-bahan}

\subsection{Perangkat Keras}
\begin{itemize}
    \item PC atau Laptop dengan arsitektur prosesor basis \textit{host} x86/x64 yang mendukung kelanjutan kompilasi emulator.
    \item Memori minimal 4GB RAM ruang eksekusi.
\end{itemize}

\subsection{Perangkat Lunak (Bahan Praktikum)}
\begin{itemize}
    \item OS \textit{Host} utama (Windows dengan WSL/GitBash, Linux, atau macOS).
    \item Text Editor (misalnya Visual Studio Code) untuk menyunting berkas kode bahasa C dan \textit{Assembly}.
    \item Instalasi Docker Engine (\texttt{docker.io}) untuk menyalakan lingkungan isolasi kompilasi.
    \item Emulator piranti keras x86: \texttt{qemu-system-x86} (atau instalasi QEMU setara) yang akan mengeksekusi binari hasil rakitan OS.
    \item Direktori \textit{Starter Code} OS 32-bit yang sudah disiapkan sebelumnya (berisi direktori \texttt{src} dengan kerangka lengkap \texttt{kernel.c}, \texttt{thread.h}, \texttt{thread.c}, \texttt{switch.asm}, \texttt{scheduler.h}, dsb).
\end{itemize}

% --- 4. Langkah-Langkah Praktikum ---
\newpage
\section{Langkah-Langkah Praktikum}

\subsection{Sesi 1: Persiapan Lingkungan Kerja}
\label{sec:m08-sesi-1}
Langkah pertama adalah menyiapkan direktori \textit{Starter Code} dan masuk ke dalam kontainer Docker.

\begin{enumerate}
    \item Buka terminal dan pastikan Anda berada di dalam folder \texttt{starter-code-os32-main}.
    \item Masuk ke lingkungan Docker dengan menjalankan perintah:
    \begin{lstlisting}[language=bash, caption={Masuk ke lingkungan kontainer}]
./run-docker.sh
    \end{lstlisting}
    \item Buka berkas \texttt{src/thread.c} pada teks editor. Berkas ini adalah file utama yang akan dimodifikasi.
\end{enumerate}

\subsection{Sesi 2: Implementasi Fungsi \texttt{thread\_create()}}
\label{sec:m08-sesi-2}
Fungsi \texttt{thread\_create()} bertugas menyiapkan \textit{Thread Control Block} (TCB) dan \textit{stack} awal untuk \textit{thread} baru. Buka \texttt{thread.c}, cari fungsi pengisian \texttt{thread\_create}, dan ikuti langkah berikut:

\begin{enumerate}
    \item \textbf{Langkah 1 (Alokasi TCB):} Gunakan \texttt{allocate\_thread\_page()} untuk alokasi memori \textit{thread}. Jika berhasil, panggil \texttt{init\_thread(t, name, priority)} dan berikan ID dengan \texttt{t->tid = allocate\_tid()}.
    \begin{figure}[H]
        \centering
        \includegraphics[width=0.8\textwidth]{Figure/08/1.png}
        \caption{Implementasi Alokasi TCB pada fungsi thread\_create()}
    \end{figure}
    
    \item \textbf{Langkah 2 (Frame 1 --- \texttt{kernel\_thread\_frame}):} Buat variabel \textit{frame} ini menggunakan \texttt{alloc\_frame()}. Isi \texttt{eip} dengan \texttt{NULL}, lalu isi \texttt{function} dan \texttt{aux} sesuai dengan argumen parameter yang dikirim.
    \begin{figure}[H]
        \centering
        \includegraphics[width=0.8\textwidth]{Figure/08/2.png}
        \caption{Implementasi Frame 1 pada fungsi thread\_create()}
    \end{figure}
    
    \item \textbf{Langkah 3 (Frame 2 --- \texttt{switch\_entry\_frame}):} Buat variabel \textit{frame} ini dengan alokasi di bawah Frame 1. Isi \texttt{eip} dengan memanggil alamat fungsi \texttt{kernel\_thread}.
    \begin{figure}[H]
        \centering
        \includegraphics[width=0.8\textwidth]{Figure/08/3.png}
        \caption{Implementasi Frame 2 pada fungsi thread\_create()}
    \end{figure}
    
    \item \textbf{Langkah 4 (Frame 3 --- \texttt{switch\_threads\_frame}):} Buat variabel \textit{frame} terbawah ini. Isi variabel \texttt{eip} dengan \texttt{switch\_entry}, lalu set nilai \texttt{ebp}, \texttt{ebx}, \texttt{esi}, dan \texttt{edi} menjadi \texttt{0}.
    \begin{figure}[H]
        \centering
        \includegraphics[width=0.8\textwidth]{Figure/08/4.png}
        \caption{Implementasi Frame 3 pada fungsi thread\_create()}
    \end{figure}
    
    \item \textbf{Langkah 5 (Penyelesaian):} Panggil \texttt{thread\_unblock(t)} agar dimasukkan ke dalam antrean \textit{ready}, kemudian lakukan \textit{return} id \textit{thread} tersebut (\texttt{t->tid}).
    \begin{figure}[H]
        \centering
        \includegraphics[width=0.8\textwidth]{Figure/08/5.png}
        \caption{Langkah Penyelesaian pada fungsi thread\_create()}
    \end{figure}
    
    \item Pastikan kerangka utuh fungsi Anda menyerupai format referensi ini:
    \begin{figure}[H]
        \centering
        \includegraphics[width=0.8\textwidth]{Figure/08/6.png}
        \caption{Fungsi utuh \texttt{thread\_create()} secara keseluruhan}
    \end{figure}
\end{enumerate}

\subsection{Sesi 3: Implementasi Cooperative \textit{Round-Robin}}
\label{sec:m08-sesi-3}
Bagian ini mengatur agar \textit{thread} dapat saling bergantian menggunakan wewenang CPU. Masih di dalam file \texttt{thread.c}:

\begin{enumerate}
    \item \textbf{Modifikasi \texttt{thread\_yield()}:} Ambil \textit{thread} yang saat ini berjalan menggunakan \texttt{thread\_current()}. Ubah statusnya menjadi \texttt{THREAD\_READY}. Masukkan \textit{thread} ini ke dalam antrean \texttt{ready\_queue} pada indeks \texttt{ready\_tail}. Geser urutan \texttt{ready\_tail}, tambahkan jumlah \texttt{ready\_count}, dan akhiri dengan memanggil fungsi \texttt{schedule()}.
    \begin{figure}[H]
        \centering
        \includegraphics[width=0.8\textwidth]{Figure/08/7.png}
        \caption{Implementasi fungsi \texttt{thread\_yield()}}
    \end{figure}

    \item \textbf{Modifikasi \texttt{next\_thread\_to\_run()}:} Fungsi ini berjalan secara metode FIFO (\textit{First-In-First-Out}). Jika antrean kosong (\texttt{ready\_count} bernilai 0), kembalikan \texttt{initial\_thread}. Jika isi antrean tersedia, ambil \textit{thread} dari urutan depan (\texttt{ready\_queue[ready\_head]}), geser letak indeks depan \texttt{ready\_head}, kurangi nilai hitung \texttt{ready\_count}, lalu \textit{return thread} yang diambil tersebut.
    \begin{figure}[H]
        \centering
        \includegraphics[width=0.8\textwidth]{Figure/08/8.png}
        \caption{Implementasi fungsi \texttt{next\_thread\_to\_run()}}
    \end{figure}
\end{enumerate}

\subsection{Sesi 4: Kompilasi dan Pengujian QEMU}
\label{sec:m08-sesi-4}
Setelah seluruh kode di \texttt{thread.c} selesai, jalankan pengujian sistem untuk membuktikan bahwa OS sudah sanggup menjalankan perpindahan waktu \textit{thread}.

\begin{enumerate}
    \item Buka kembali terminal yang berada di dalam lingkungan Docker, lalu jalankan perintah kompilasi:
    \begin{lstlisting}[language=bash, caption={Perintah build pada terminal Docker}]
make
    \end{lstlisting}
    \item Setelah proses kompilasi selesai (\textit{error} tidak muncul), jalankan QEMU emulator secara langsung menggunakan perintah di bawah ini:
    \begin{lstlisting}[language=bash, caption={Menjalankan OS pada terminal Docker}]
make run
    \end{lstlisting}
    \item Perhatikan layar simulator QEMU. Anda harus masuk dalam putaran sistem \textit{infinite loop} yang memperlihatkan ketiga \textit{thread} yang dibuat mencetak teks (Thread A, Thread B, dan Thread C) bergantian per waktu putaran:
    \begin{lstlisting}[caption={Ekspektasi output sistem penjadwalan Thread di latar QEMU}]
===================================
  IF-OS  |  Praktikum Sistem Operasi
===================================

Inisialisasi sistem thread...
Membuat thread A, B, C...
Memulai scheduling...

[Thread A] tick 0
[Thread B] tick 0
[Thread C] tick 0
[Thread A] tick 1
[Thread B] tick 1
...
    \end{lstlisting}
\end{enumerate}
